\documentclass[11pt,a4paper]{article}
\usepackage[utf8]{inputenc}
\usepackage{amsmath,amssymb}
\usepackage{hyperref}

\title{Redes de Hopfield e Memória Associativa Auto-organizável}
\author{John J. Hopfield \and Pesquisadores Associados}
\date{1982}

\begin{document}

\maketitle

\begin{abstract}
Este trabalho formaliza a dinâmica de sistemas neurais com conexões simétricas recorrentes e formulação energética baseada no modelo de Ising da mecânica estatística. Demonstramos que o sistema converge monotonicamente para mínimos locais de uma função de energia de Liapunov.
\end{abstract}

\section{Introdução}
Modelos conexionistas recorrentes possuem capacidade inerente de recuperação de padrões ruidosos através de dinâmica de atratores. A energia do sistema é dada por:

\begin{equation}
E = -\frac{1}{2} \sum_{i=1}^{N} \sum_{j=1}^{N} w_{ij} s_i s_j - \sum_{i=1}^{N} \theta_i s_i
\end{equation}

onde $w_{ij}$ representa o peso sináptico simétrico entre os neurônios $i$ e $j$ ($w_{ij} = w_{ji}$, com $w_{ii} = 0$), $s_i \in \{-1, +1\}$ é o estado binário de ativação, e $\theta_i$ é o limiar de disparo.

\section{Regra de Aprendizado Hebbiano}
Para armazenar $P$ padrões ortogonais $\xi^{\mu}$, os pesos são definidos pela regra hebbiana clássica:

\begin{equation}
w_{ij} = \frac{1}{N} \sum_{\mu=1}^{P} \xi_i^{\mu} \xi_j^{\mu}
\end{equation}

A capacidade teórica de armazenamento estável sem saturação é estimada em $P_{max} \approx 0.14 N$.

\section{Conclusões}
A formulação energética estabelece uma ponte entre física estatística e inteligência artificial computacional.

\end{document}
